% APORTE TERCER AVANCE (CIERRE) - GILCES CARRANZA JOSE IGNACIO (CI: 1251059240)
% Incluye: Revisión integral como Analista Líder, Conclusiones 1 y 2 (>80 palabras c/u), Declaración de Uso Ético de IA y Compilación Final

\section{Paso 6: Conclusiones Individuales y Declaración de Uso de IA}

\subsection{Conclusiones de José Ignacio Gilces Carranza}
\begin{enumerate}
	
	\item \textbf{Conclusión 1 --- Gilces Carranza José Ignacio (Sobre la Especificación EARS y la Propiedad ``Apropiado al Nivel'' en IEEE 29148):} \\
	\textbf{(1) Afirmación:} La redacción de requisitos mediante patrones sintácticos estructurados y la abstención estricta de incluir tecnologías de implementación en la especificación eliminan la ambigüedad semántica y preservan la independencia entre el problema de negocio y el diseño posterior. \textbf{(2) Evidencia con cita:} Conforme lo establecen la norma ISO/IEC/IEEE 29148:2018 (cláusula 5.2.5) \cite{ieee29148} y la metodología EARS de Mavin et al. \cite{mavin2009ears}, todo requisito individual bien formado debe ser inambiguo, singular, verificable y apropiado al nivel (describiendo el ``qué'' necesita el stakeholder sin imponer prematuramente el ``cómo'' tecnológico). \textbf{(3) Implicación para el PFC y la profesión:} En nuestra Plataforma de Ventas para almacenes eléctricos de Quevedo, depurar los 15 requisitos de menciones a herramientas de programación específicas y expresarlos mediante umbrales observables ($\le 3.0\text{ s}$ en recomendaciones técnicas filtradas por existencias reales $>0$ en el catálogo local) nos permitió construir una especificación limpia, verificable y profesionalmente auditable.
	
	\item \textbf{Conclusión 2 --- Gilces Carranza José Ignacio (Sobre la Cuantificación de Atributos de Calidad con ISO/IEC 25010:2023):} \\
	\textbf{(1) Afirmación:} La viabilidad operativa de un sistema transaccional acoplado a asistencia inteligente depende decisivamente de la cuantificación rigurosa de sus Requisitos No Funcionales (RNF) en términos de eficiencia de desempeño, fiabilidad, seguridad y capacidad de interacción. \textbf{(2) Evidencia con cita:} De acuerdo con el estándar ISO/IEC 25010:2023 \cite{iso25010} y la clasificación de Glinz \cite{glinz2007nfr}, los atributos de calidad acotan el comportamiento del producto bajo condiciones reales de operación y definen criterios objetivos de aceptación. \textbf{(3) Implicación para el PFC y la profesión:} Vincular cada RNF del PFC al modelo ISO/IEC 25010 nos permitió exigir una disponibilidad del $\ge 99.5\%$, un mecanismo de contingencia local en $\le 200\text{ ms}$ ante interrupciones del servicio externo de recomendación de IA, sesiones autenticadas por turno con cifrado irreversible y estándares de accesibilidad visual WCAG 2.1 AA con operación íntegra por teclado, garantizando continuidad absoluta en el cobro en horas pico.
	
\end{enumerate}

\subsection{Declaración Explícita de Uso Ético y Validación Crítica de Inteligencia Artificial}
En estricto cumplimiento de la Sección 4 y del Paso 6d de la Guía de Práctica Experimental de la UTEQ, los autores (\textbf{Gilces Carranza José Ignacio, Vélez Vinces Jostin Gustavo y Robinson Espinoza Jeanpierre}) declaran de manera transparente que durante la elaboración del presente informe se utilizaron herramientas de Inteligencia Artificial Generativa como apoyo de asistencia en la estructuración tipográfica, revisión de estilo técnico y verificación sintáctica de plantillas EARS. Todo el contenido conceptual, la selección y contextualización del sistema real en Quevedo (PFC), las métricas cuantitativas de los 15 requisitos (verificando que se mantengan libres de decisiones prematuras de implementación), el análisis de las ocho perspectivas de contexto y las citas bibliográficas fueron contrastados, verificados y validados críticamente por los integrantes del equipo contra las fuentes primarias obligatorias (ISO/IEC/IEEE 29148:2018 \cite{ieee29148}, ISO/IEC 25010:2023 \cite{iso25010}, Pohl 2025 \cite{pohl2025} y SWEBOK v4.0 \cite{swebok2024}), asumiendo los autores la responsabilidad académica e intelectual íntegra sobre el documento final.